\subsection{相伴纤维丛}

设 E 与 F 为集合，G 为作用于 E 的右边、作用于 F 的左边的群。群 G 作用于乘积 \(E \times F\) 的右边，其法则为 \((x, y) \cdot g = (x \cdot g, g^{-1} \cdot y)\) ，其中 \(g \in G\)，\((x, y) \in E \times F\) 。 \(E \times F\) 对此作用的商集记作 \(E \times^{G} F\) 。

当 E 与 F 是拓扑空间时，给集合 \(E \times^{G} F\) 赋予由 \(E \times F\) 的拓扑诱导的商拓扑。从 \(E \times F\) 到 \(E \times^{G} F\) 上的典范映射是连续的。若群 G 连续地作用于 E 与 F，则它是开的。

另外，设 \(F'\) 为 G 从左边作用于其上的集合，\(h: F \to F'\) 为与 G 在 F 与 \(F'\) 上的作用相容的映射（A, I, 第 50 页）。映射 \(Id_{E} \times h: E \times F \to E \times F'\) 与 G 的作用相容，并且通过过渡到商，定义一个记作 \(Id_{E} \times^{G} h\) 的、从 \(E \times^{G} F\) 到 \(E \times^{G} F'\) 中的映射。

若 \(h \colon F \to F'\) 是连续映射（与 G 在 F 与 \(F'\) 上的作用相容），则映射 \(\mathrm{Id}_{\mathrm{E}} \times^{\mathrm{G}} h\) 是连续的（TG, I, 第 21 页, 命题 6）。

例. —— 1) 设 F 为拓扑空间，G 为从左边连续作用于 F 的拓扑群。若给拓扑空间 G 赋予 G 以右平移的作用，则空间 G × \(^{G}\) F 典范地等同于 F，如下所述。连续映射 \(\varphi\) : F → G × F 与 \(\psi\) : G × F → F （分别由 \(\varphi(f) = (e, f)\) （其中 e 表示 G 的单位元）与 \(\psi(g, f) = g \cdot f\) 定义）诱导出连续映射 \(\overline{\varphi}\) : F → G × \(^{G}\) F 与 \(\overline{\psi}\) : G × \(^{G}\) F → F，二者互为逆映射。

\begin{enumerate}
\def\labelenumi{\arabic{enumi})}
\setcounter{enumi}{1}
\item
  例 1 可如下推广。设 B 与 F 为拓扑空间，G 为从左边连续作用于 F 的拓扑群。通过过渡到商，由 (b, f) ↦ ((b, e), f) 给出的从 B × F 到 (B × G) × F 的映射，以及由 ((b, g), f) ↦ (b, gf) 给出的从 (B × G) × F 到 B × F 的映射，定义出互为逆映射的 B-同构 B × F → (B × G) × \(^{G}\) F 与 (B × G) × \(^{G}\) F → B × F。
\item
  类似地，设 E 为拓扑空间，G 为从右边连续作用于 E 的拓扑群。若给拓扑空间 G 赋予 G 以左平移的作用，则空间 \(E \times^{G} G\) 等同于 E。
\end{enumerate}

设 E 与 F 为拓扑空间，G 为从右边连续作用于 E、从左边连续作用于 F 的拓扑群。设 B 为拓扑空间，\(p\colon E \to B\) 为满足 \(p(x \cdot g) = p(x)\) （其中 \(x \in E\)，\(g \in G\)）的连续映射。映射 \(p \circ pr_{1}\colon E \times F \to B\) 通过过渡到商，定义一个连续映射 \(p^{F}\colon E \times^{G} F \to B\) ，而典范映射 \(\pi\colon E \times F \to E \times^{G} F\) 是一个 B-态射。

设 B' 为拓扑空间，\(h\colon B'\to B\) 为连续映射。群 G 连续地作用于 \(B'\times_{B}E\) 的右边，其作用法则为 \(((b',x),g)\mapsto(b',x\cdot g)\)。将典范同构 \(((b',x),y)\mapsto(b',(x,y))\) ——它把 \((\mathrm{B}'\times_{\mathrm{B}}\mathrm{E})\times\mathrm{F}\) 映到 \(\mathrm{B}'\times_{\mathrm{B}}(\mathrm{E}\times\mathrm{F})\) 上——与映射 \(\operatorname{Id}_{\mathrm{B}'}\times\pi\) （从 \(\mathrm{B}'\times_{\mathrm{B}}(\mathrm{E}\times\mathrm{F})\) 到 \(\mathrm{B}'\times_{\mathrm{B}}(\mathrm{E}\times^{\mathrm{G}}\mathrm{F})\) 中）复合，得到连续映射 \(\lambda_{0}\colon(\mathrm{B}'\times_{\mathrm{B}}\mathrm{E})\times\mathrm{F}\to\mathrm{B}'\times_{\mathrm{B}}(\mathrm{E}\times^{\mathrm{G}}\mathrm{F})\)。它通过过渡到商，定义一个连续映射

\[
\lambda \colon (\mathrm{B} ^ {\prime} \times_ {\mathrm{B}} \mathrm{E}) \times^ {\mathrm{G}} \mathrm{F} \rightarrow \mathrm{B} ^ {\prime} \times_ {\mathrm{B}} (\mathrm{E} \times^ {\mathrm{G}} \mathrm{F}).
\]

引理. —— 映射 λ 是同胚。

映射 \(\lambda_0\) 是满射，且 \((\mathrm{B}' \times_{\mathrm{B}} \mathrm{E}) \times \mathrm{F}\) 中两个元素在 \(\lambda_0\) 下有同一像，当且仅当它们属于 G 在 \((\mathrm{B}' \times_{\mathrm{B}} \mathrm{E}) \times \mathrm{F}\) 上作用的同一轨道。由此可知映射 \(\lambda\) 是双射。由于映射 \(\pi\) 是开的，映射 \(\mathrm{Id}_{\mathrm{B}'} \times_{\mathrm{B}} \pi\) 也是开的（I, 第 17 页, 命题 8），从而 \(\lambda_0\) 与 \(\lambda\) 也是开的（TG, I, 第 32 页, 命题 3），这就证明了 \(\lambda\) 是同胚。

命题 7. —— 设 B 为拓扑空间，G 为拓扑群，(E,p) 为 B 上以 G 为群的主纤维丛，F 为 G 从左边连续作用于其上的拓扑空间。

\begin{enumerate}
\def\labelenumi{\alph{enumi})}
\item
  拓扑空间 \(E \times^{G} F\) 连同连续映射 \(p^{F}: E \times^{G} F \to B\) （由映射 \(p \circ pr_{1}: E \times F \to B\) 过渡到商而得），是以 F 为纤维型的局部平凡纤维丛；若 B-丛 E 可平凡化，则它也可平凡化。
\item
  设 \(\pi\colon\mathrm{E}\times\mathrm{F}\to\mathrm{E}\times^{\mathrm{G}}\mathrm{F}\) 为典范满射。映射 \(\mu\colon\mathrm{E}\times\mathrm{F}\to\mathrm{E}\times_{\mathrm{B}}(\mathrm{E}\times^{\mathrm{G}}\mathrm{F})\) 使 \((x,\pi(x,f))\) 对应于 \((x,f)\) ，它是一个同胚，其逆映射是 E 上的局部平凡纤维丛 \((\mathrm{E}\times_{\mathrm{B}}(\mathrm{E}\times^{\mathrm{G}}\mathrm{F}),\mathrm{pr}_{1})\) 的一个平凡化。
\end{enumerate}

先设 E 为平凡主纤维丛 \(B \times G\) 。由例 2，此时空间 \(E \times^{G} F\) 等同于 \(B \times F\) ，映射 \(p^{F}\) 等同于第一投影 \(pr_{1}: B \times F \to B\) ，这就证明在此情形 \((\mathrm{E} \times^{\mathrm{G}}\mathrm{F}, p^{\mathrm{F}})\) 是 B 上的可平凡化纤维丛。在一般情形，B 的每个点都有邻域 U，使映射 \(p_{\mathrm{U}}: p^{-1}(\mathrm{U}) \to \mathrm{U}\) 把 \(p^{-1}(\mathrm{U})\) 变为 U 上的可平凡化主纤维丛。由前所述，\((p^{-1}(\mathrm{U}) \times^{\mathrm{G}}\mathrm{F} \to \mathrm{U}, (p_{\mathrm{U}})^{\mathrm{F}})\) 是 U 上的可平凡化纤维丛。把上面的引理应用于 \(B' = U\) 且 \(h: U \to B\) 为典范单射的情形，对 U-空间 \(((p^{\mathrm{F}})^{-1}(\mathrm{U}), (p^{\mathrm{F}})_{\mathrm{U}})\) （由 \((\mathrm{E} \times^{\mathrm{G}}\mathrm{F}, p^{\mathrm{F}})\) 限制到 U 上得到）结论同样成立。这就证明 \((\mathrm{E} \times^{\mathrm{G}}\mathrm{F}, p^{\mathrm{F}})\) 是 B 上的局部平凡纤维丛，从而完成了断言 a) 的证明。

映射 \(\theta\colon\mathrm{E}\times\mathrm{G}\to\mathrm{E}\times_{\mathrm{B}}\mathrm{E}\) 由 \(\theta(x,g)=(x,x\cdot g)\) 所定义，它是同胚（I, 第 94 页, 命题 3），且与 G 在 E \(\times\) G 上和 E \(\times_{\mathrm{B}}\) E 上由 \(((x,g),g')\mapsto(x,gg')\) 与 \(((x,y),g')\mapsto(x,yg')\) 给出的作用相容。因此 \(\theta\) 通过过渡到商，诱导出同胚 \(\theta'\) （从 (E \(\times\) G) \(\times^{\mathrm{G}}\) F 到 (E \(\times_{\mathrm{B}}\) E) \(\times^{\mathrm{G}}\) F 上）。映射 \(\mu\) 是同胚 E \(\times\) F \(\rightarrow\) (E \(\times\) G) \(\times^{\mathrm{G}}\) F （例 2）、\(\theta'\) 以及典范同胚 (E \(\times_{\mathrm{B}}\) E) \(\times^{\mathrm{G}}\) F \(\rightarrow\) E \(\times_{\mathrm{B}}\) (E \(\times^{\mathrm{G}}\) F) （I, 第 105 页, 引理）的复合，由此即得 \(b\) 。

在命题 7 的假设之下，局部平凡纤维丛 (E × \(^{G}\) F, p \(^{F}\) ) 称为与主纤维丛 (E, p) 相伴的、以 F 为纤维型的局部平凡纤维丛。B-空间 E × \(^{G}\) F 的一切纤维都与空间 F 同胚。特别地，若空间 F 是离散的，则映射 p \(^{F}\) 是一个覆叠。

例. —— 4) 设 B 为拓扑空间，G 为拓扑群，(E, p) 为 B 上的主 G-丛。设 H 为 G 的子群。

以 \(\varphi\colon\mathrm{E}\to\mathrm{E}\times(\mathrm{G}/\mathrm{H})\) 与 \(\psi\colon\mathrm{E}\times(\mathrm{G}/\mathrm{H})\to\mathrm{E}/\mathrm{H}\) 表示由 \(\varphi(x)=(x,\mathrm{H})\) 与 \(\psi(x,g\mathrm{H})=(x\cdot g)\mathrm{H}\) 定义的映射。它们与到 B 的投影以及 G 的作用相容，并通过过渡到商，定义出互为逆映射的 B-空间态射 \(\overline{\varphi}\colon\mathrm{E}/\mathrm{H}\to\mathrm{E}\times^{\mathrm{G}}(\mathrm{G}/\mathrm{H})\) 与 \(\overline{\psi}\colon\mathrm{E}\times^{\mathrm{G}}(\mathrm{G}/\mathrm{H})\to\mathrm{E}/\mathrm{H}\) 。我们称 \(\overline{\varphi}\) 为从 E/H 到 \(\mathrm{E}\times^{\mathrm{G}}(\mathrm{G}/\mathrm{H})\) 上的典范同胚。特别地，拓扑空间 E/H 连同连续映射 \(p_{\mathrm{H}}\colon\mathrm{E}/\mathrm{H}\to\mathrm{B}\) ，是 B 上的以 G/H 为纤维型的局部平凡纤维丛。

此外，若 H 是 G 的正规子群，则 G 的作用通过过渡到商，赋予 B-空间 E/H 一个以 G/H 为群的主纤维丛结构。

特别地，若 E 是 B 的以 G 为群的主覆叠，则 E/H 是 B 的覆叠；若 H 在 G 中正规，则它是以 G/H 为群的主覆叠。

\begin{enumerate}
\def\labelenumi{\arabic{enumi})}
\setcounter{enumi}{4}
\item
  设 B 为拓扑空间，G 为拓扑群，E 为 B 上的主 G-丛。设 F 为关于 G 的拓扑齐性空间（TG, III, 第 12 页）。设 y 为 F 的一点，\(G_{y}\) 为其稳定子。映射 \(\varphi_{y}\colon x\mapsto(x,y)\) （从 E 到 \(E\times F\)）通过过渡到商，定义出同胚 \(\overline{\varphi}_{y}\) （从 \(E/G_{y}\) 到 \(E\times^{G}F\) 上）。当群 G 交换时，子群 \(G_{y}\) 不依赖于点 y，但同胚 \(\overline{\varphi}_{y}\) 一般依赖于它。
\item
  设 B 为拓扑空间，G 为拓扑群，\((E,p)\) 为 B 上的主 G-丛。设 H 为拓扑群，\(f:G\to H\) 为拓扑群态射；给 H 赋予由 \(g\cdot h=f(g)h\) 给出的 G 的左作用。
\end{enumerate}

设 \(q \colon \mathrm{E} \times \mathrm{H} \to \mathrm{E} \times^{\mathrm{G}} \mathrm{H}\) 为典范映射；它是开的，从而是普遍严格的（I, 第 20 页, 推论 11）。映射 \(m \colon (x, h, h') \mapsto q(x, hh')\) （从 \(\mathrm{E} \times \mathrm{H} \times \mathrm{H}\) 到 \(\mathrm{E} \times^{\mathrm{G}} \mathrm{H}\) 中）是连续的。设 \((x, h) \in \mathrm{E} \times \mathrm{H}\) ，\(h' \in \mathrm{H}\) 以及 \(g \in \mathrm{G}\) ；我们有 \(m(xg, f(g)^{-1}h, h') = q(xg, f(g)^{-1}hh') = q(x, hh') = m(x, h, h')\) 。因此，存在唯一的映射

\[
m ^ {\prime} \colon (\mathrm{E} \times^ {\mathrm{G}} \mathrm{H}) \times \mathrm{H} \rightarrow \mathrm{E} \times^ {\mathrm{G}} \mathrm{H}
\]

使得 \(m'(q(x,h),h') = q(x,hh')\) 对一切 \(x \in E\) 与一切 \(h, h' \in H\) 成立。由于 q 是普遍严格的，映射 \(m'\) 是连续的。这是拓扑群 H 在 B-空间 \(E \times^{G} H\) 上的一个右作用。

设 U 为 B 的开子集，使得 \(E_{U}\) 同构于 U 上的主纤维丛 \(U \times G\) 。赋予 H 的作用后，U-空间 \((\mathrm{E} \times^{\mathrm{G}}\mathrm{H})_{\mathrm{U}}\) 等同于主纤维丛 \(U \times H\) 。这就证明 \(E \times^{G}H\) 是 B 上的主 H-丛。

定义 4. —— 设 B 为拓扑空间，G 为拓扑群。称 B 上的局部平凡纤维丛 X 与 B 上以 G 为群的主纤维丛 E 相伴，如果存在一个 G 从左边连续作用于其上的拓扑空间 F，以及一个从 E × \(^{G}\) F 到 X 上的 B-同构。

设 E 为 B 上的主 G-丛，X 为 B 上与 E 相伴的局部平凡纤维丛。若主丛 E 可平凡化，则 X 可平凡化（命题 7）。若 B' 为拓扑空间且 \(h \colon \mathrm{B}' \to \mathrm{B}\) 为连续映射，则由 X 经换基得到的局部平凡纤维丛 \(\mathrm{B}' \times_{\mathrm{B}} \mathrm{X}\) 与 B' 上的主纤维丛 \(\mathrm{B}' \times_{\mathrm{B}} \mathrm{E}\) 相伴（I, 第 105 页, 引理）。

